% Преамбула отчёта (pandoc + XeTeX): заголовки PT Sans, подписи по-русски, нумерация по разделам,
% мелкий шрифт таблиц, альбомные страницы, рисунки на месте, без висящих заголовков, приложения А–Д.
\usepackage{etoolbox}
\usepackage{float}
\usepackage{pdflscape}
\usepackage{needspace}
\usepackage[font=small,labelfont=bf,labelsep=endash,justification=centering,singlelinecheck=true]{caption}
\usepackage{titlesec}
\usepackage{xurl}
\usepackage{newunicodechar}
\makeatletter
\def\fps@figure{tp}
\makeatother
\usepackage[section]{placeins}
\renewcommand{\topfraction}{0.9}\renewcommand{\bottomfraction}{0.8}\renewcommand{\textfraction}{0.07}
\renewcommand{\floatpagefraction}{0.8}
\counterwithin{figure}{section}
\counterwithin{table}{section}
\counterwithin{equation}{section}
\AtBeginDocument{%
  \renewcommand{\figurename}{Рисунок}%
  \renewcommand{\tablename}{Таблица}%
  \renewcommand{\contentsname}{Содержание}}
% заголовки: раздел — с новой страницы; в приложениях — «Приложение А.»
\newif\ifinappendix
\titleformat{\section}{\normalfont\sffamily\Large\bfseries}{\ifinappendix Приложение~\thesection.\else\thesection\fi}{0.6em}{}
\titleformat{\subsection}{\normalfont\sffamily\large\bfseries}{\thesubsection}{0.6em}{}
\titleformat{\subsubsection}{\normalfont\sffamily\normalsize\bfseries}{\thesubsubsection}{0.6em}{}
\titlespacing*{\section}{0pt}{4ex plus 1ex minus .3ex}{1.4ex plus .3ex}
\titlespacing*{\subsection}{0pt}{2.6ex plus 1ex minus .2ex}{1ex plus .2ex}
\titlespacing*{\subsubsection}{0pt}{2ex plus .8ex minus .2ex}{0.8ex plus .2ex}
\pretocmd{\section}{\needspace{14\baselineskip}}{}{}
% титул и оглавление — на отдельных страницах
\AtBeginDocument{\apptocmd{\maketitle}{\thispagestyle{empty}\clearpage}{}{}%
  \apptocmd{\tableofcontents}{\clearpage}{}{}}
\pretocmd{\subsection}{\needspace{11\baselineskip}}{}{}
\pretocmd{\subsubsection}{\needspace{10\baselineskip}}{}{}
% приложения: буквы А–Е, уникальные якоря hyperref
\newcommand{\startappendix}{%
  \clearpage\appendix\inappendixtrue
  \renewcommand{\thesection}{\ifcase\value{section}\or А\or Б\or В\or Г\or Д\or Е\fi}%
  \renewcommand{\theHsection}{app\arabic{section}}}
% таблицы — мелким шрифтом и плотнее
\AtBeginEnvironment{longtable}{\footnotesize\setlength{\tabcolsep}{4pt}\renewcommand{\arraystretch}{1.15}}
\setlength{\emergencystretch}{3em}
\widowpenalty=10000
\clubpenalty=10000
% символы, которых нет в PT Serif / PT Sans / PT Mono, — из STIXGeneral
\newfontfamily\fallbackfont{STIXGeneral}
\newunicodechar{τ}{{\fallbackfont\upshape τ}}
\newunicodechar{Σ}{{\fallbackfont\upshape Σ}}
\newunicodechar{α}{{\fallbackfont\upshape α}}
\newunicodechar{→}{{\fallbackfont\upshape →}}
\newunicodechar{∈}{{\fallbackfont\upshape ∈}}
\newunicodechar{⊂}{{\fallbackfont\upshape ⊂}}
% уникальные якоря hyperref при нумерации по разделам
\AtBeginDocument{%
  \renewcommand{\theHtable}{\theHsection.\arabic{table}}%
  \renewcommand{\theHfigure}{\theHsection.\arabic{figure}}%
  \renewcommand{\theHequation}{\theHsection.\arabic{equation}}}
% блоки кода: мельче и с переносом длинных строк
\usepackage{fvextra}
\AtBeginDocument{\RecustomVerbatimEnvironment{Highlighting}{Verbatim}{commandchars=\\\{\},breaklines,breakanywhere,fontsize=\footnotesize}}
